\section{Order Sealing and Transport}
\label{sec:orders}

\begin{figure}[t]
\centering
\begin{tikzpicture}[
  font=\small,
  head/.style={draw, rounded corners=2pt, fill=black!5, align=center, minimum width=2.3cm, minimum height=0.8cm},
  msg/.style={-{Stealth[length=2mm]}, thick},
  note/.style={font=\scriptsize, align=left, fill=white, inner sep=1.5pt},
]
\def\xc{0} \def\xb{3.5} \def\xv{8.0} \def\xa{11.5}
\node[head] (hc) at (\xc,0) {Buyer client};
\node[head] (hb) at (\xb,0) {Buyer node};
\node[head] (hv) at (\xv,0) {Vendor node\\\scriptsize mailbox};
\node[head] (ha) at (\xa,0) {Vendor client};
\foreach \x in {\xc,\xb,\xv,\xa} { \draw[black!40] (\x,-0.45) -- (\x,-9.2); }

\draw[msg] (\xb,-1.0) -- node[note, above]{1. verified catalog} (\xc,-1.0);
\draw[msg] (\xc,-1.9) -- node[note, above]{2. checkout} (\xb,-1.9);
\node[note, right] at (\xb+0.1,-2.7) {3. seal to the seller's mailbox key};
\draw[msg] (\xb,-3.5) -- node[note, above]{4. letter, retried from an outbox} (\xv,-3.5);
\node[note, left] at (\xv-0.1,-4.25) {5. price the order;\\\phantom{5. }fresh address};
\draw[msg] (\xv,-5.15) -- node[note, above]{6. terms, or a decline} (\xb,-5.15);
\node[note, right] at (\xb+0.1,-5.9) {7. pay, and name the transaction};
\node[note, left] at (\xv-0.1,-6.6) {8. ``payment sent'' until a block};
\draw[msg] (\xa,-7.4) -- node[note, above]{9. ship once paid} (\xv,-7.4);
\draw[msg] (\xb,-8.3) -- node[note, above]{10. buyer confirms} (\xv,-8.3);
\end{tikzpicture}
\caption{A sale between two nodes (\S\ref{sec:sale}). Steps 4, 6, and 10 are sealed letters to the other node's loopback mail endpoint. The seller shows step 8 because compact filters cannot see the mempool. A multi-seller checkout repeats the exchange once per seller.}
\label{fig:orderflow}
\end{figure}


An order is a private letter to one vendor, not a catalog record and not an event on a public relay. Figure~\ref{fig:orderflow} shows the flow.

\subsection{Order contents}

The order plaintext (magic \code{OBOD}, version 1) contains: the address of the first listing; its quantity; an options map of up to 16 string pairs, which carries the settlement type, the shipping charge, a delivery note for digital goods, and the address and quantity of any further listings from the same seller; the ship-to text (at most 1024 bytes) for physical goods; the payment rail; and a fresh X25519 public key $R$ that the vendor uses to reply. The order carries no long-term buyer identity key and no buyer signature. The vendor learns who the buyer is only from what the buyer writes into the letter (for a physical good, the ship-to) and from the payment.

\subsection{Sealing}

The node seals the plaintext with HPKE~\citep{rfc9180} in base mode with the suite DHKEM(X25519, HKDF-SHA256)~\citep{rfc7748,rfc5869}, HKDF-SHA256, and ChaCha20-Poly1305~\citep{rfc8439}. The recipient key is the vendor's mailbox key $P_V$; the HPKE \code{info} string is \code{openbazaar-tormail-order-v1}; the associated data is empty. The sealed letter is
\[
  \code{OBSL} \,\|\, \mathtt{0x01} \,\|\, \mathit{enc} \,\|\, \mathrm{len}_{32}(\mathit{ct}) \,\|\, \mathit{ct},
\]
where $\mathit{enc}$ is HPKE's 32-byte encapsulated ephemeral key. Letters larger than 64~KiB are rejected. A vendor reply is sealed the same way to the buyer's per-order key $R$ with \code{info} \code{openbazaar-tormail-reply-v1} and a size limit of 8~KiB. Because $R$ is generated per order and travels inside the ciphertext, no relay or mailbox needs to remember which buyer wrote to which vendor in order to route a reply.

Base mode means letters are not authenticated to a sender: anyone who knows $P_V$ can produce a letter that the vendor will decrypt. This is deliberate, since a buyer should not need a long-term identity to order, but it means the vendor must treat an order as an offer that becomes real when paid. It also means that a vendor's mailbox can be filled with junk letters; the mailbox caps letter size but does not rate-limit.

Letters are not padded. The ciphertext is exactly 16 bytes longer than the plaintext, so an observer who sees a letter at rest or in transit learns its length, which correlates with the length of the ship-to and the number of lines. Over Tor, a letter of a few hundred bytes fits in a small number of fixed-size cells, which hides the exact length from relays but not the approximate size. Padding to fixed size buckets would hide it (\S\ref{sec:future}).

\subsection{Multi-seller checkout}

A checkout may contain packages from several sellers. The node validates every package (each listing exists in the verified catalog, each package is one seller and each seller one package, moderated packages only contain listings that accept the moderator) before sealing anything, then seals one letter per seller to that seller's key. A seller learns only the lines in its own package. Each package gets its own payment instructions and its own payment (\S\ref{sec:payments}).

\subsection{The mailbox}

The vendor's node accepts letters into a mailbox directory without loading any secret: a letter is written to a temporary file and renamed to a random 128-bit identifier, so the identifier is not derived from the sender. Letters are opened later with the mailbox secret. A node that is compromised while the vendor is away therefore holds ciphertext, not orders, unless the mailbox secret is also stored on it.

\subsection{Transport}

Letters travel over Tor~\citep{dingledine2004tor}. The vendor's node runs a v3 onion service~\citep{torrendv3} whose address is published with its store record. The buyer's node dials it on a fixed virtual port and writes one framed letter (\code{OBML}, version 1); the listener accepts one stream per rendezvous, rejects streams for any other port, and writes the letter into the mailbox. If the vendor's node is unreachable, the buyer's node keeps the letter in an outbox and retries until a deadline, after which the buyer is told the order was not delivered. No third-party relay stores letters: a vendor whose node is never awake is a closed shop. Escrow letters (\S\ref{sec:escrow}) and dispute letters use the same transport to the counterparty's or moderator's onion mailbox.

Using an onion service for the vendor's mailbox means the vendor's IP address is not revealed to buyers, and the buyer's node connects through a Tor circuit, so the vendor does not learn the buyer's IP address. Tor does not protect against an adversary who observes both ends of a circuit~\citep{dingledine2004tor,johnson2013users}; \S\ref{sec:privacy} discusses this. Carrying letters over a mix network such as Nym~\citep{diaz2021nym,piotrowska2017loopix} to resist timing correlation is future work (\S\ref{sec:future}).
